% ============================================================================
% CHAPTER 2: STATISTICAL NLP AND PROBABILISTIC MODELS
% ============================================================================
% Covers n-gram language models, perplexity, HMMs, CRFs, topic models,
% and statistical machine translation.  The probabilistic backbone of NLP.
% ============================================================================

\chapter{Statistical NLP and Probabilistic Models}
\label{chap:statistical_nlp}

\begin{tldr}
Statistical NLP models---n-gram language models, Hidden Markov Models, and
Conditional Random Fields---form the probabilistic backbone of natural language
processing.  These are not historical curiosities: CRF layers still appear on
top of BERT for NER, perplexity remains the primary language model evaluation
metric, and understanding these models demonstrates the probabilistic thinking
that interviewers value.  This chapter covers language modeling, sequence
labeling, and topic models with the mathematical rigor expected at L5+ levels.
If you can derive the forward algorithm, explain Kneser-Ney smoothing, and
articulate why CRFs capture label dependencies that softmax classifiers miss,
you signal genuine depth that sets you apart from candidates who only know the
API layer.
\end{tldr}

\bigskip

% ============================================================================
% SECTION 1: N-GRAM LANGUAGE MODELS
% ============================================================================
\section{N-gram Language Models}
\label{sec:ngram_lm}

A \textbf{language model} assigns a probability to a sequence of words.  Given
a sentence $W = w_1, w_2, \ldots, w_N$, a language model computes
$P(w_1, w_2, \ldots, w_N)$.  This is the single most fundamental concept in
NLP: every spell-checker, speech recognizer, machine translation system, and
large language model is, at its core, a language model.

\subsection{The Chain Rule Decomposition}

By the chain rule of probability, any joint distribution over a sequence can
be decomposed exactly as:

\begin{mathresult}{Chain Rule for Language Modeling}
\begin{equation}
P(w_1, w_2, \ldots, w_N) = \prod_{t=1}^{N} P(w_t \given w_1, w_2, \ldots, w_{t-1})
\label{eq:chain_rule}
\end{equation}
This is \emph{exact}---no assumptions have been made.  The entire history of
language modeling is about how to estimate each conditional factor
$P(w_t \given w_1, \ldots, w_{t-1})$.
\end{mathresult}

The problem is immediately apparent: estimating $P(w_t \given w_1, \ldots, w_{t-1})$
requires conditioning on the \emph{entire} preceding context, which grows
without bound.  For a vocabulary of size $V$ and context length $t-1$, there
are $V^{t-1}$ possible contexts---a number that explodes far beyond any
dataset size.

\subsection{The Markov Assumption}

The n-gram approach makes a simplifying \textbf{Markov assumption}: the
probability of the next word depends only on the preceding $n-1$ words, not
the full history.

\begin{mathresult}{N-gram Markov Assumption}
\begin{equation}
P(w_t \given w_1, \ldots, w_{t-1}) \approx P(w_t \given w_{t-n+1}, \ldots, w_{t-1})
\label{eq:markov}
\end{equation}
For a bigram model ($n=2$): $P(w_t \given w_1, \ldots, w_{t-1}) \approx P(w_t \given w_{t-1})$.\\
For a trigram model ($n=3$): $P(w_t \given w_1, \ldots, w_{t-1}) \approx P(w_t \given w_{t-2}, w_{t-1})$.
\end{mathresult}

The tradeoff is fundamental: larger $n$ captures more context but requires
exponentially more data to estimate ($V^n$ possible n-grams).  In practice,
$n = 3$ to $5$ was typical for statistical systems, with trigram models being
the dominant workhorse for decades.

\subsection{Maximum Likelihood Estimation}

N-gram probabilities are estimated from corpus counts using maximum likelihood
estimation (MLE):

\begin{mathresult}{N-gram MLE}
\begin{equation}
P_{\text{MLE}}(w_t \given w_{t-n+1}, \ldots, w_{t-1})
= \frac{C(w_{t-n+1}, \ldots, w_t)}{C(w_{t-n+1}, \ldots, w_{t-1})}
\label{eq:ngram_mle}
\end{equation}
where $C(\cdot)$ denotes the count of an n-gram or $(n{-}1)$-gram in the
training corpus.
\end{mathresult}

\paragraph{Bigram example.}
Consider a training corpus: ``I like cats. I like dogs. I hate rain.''
\begin{align*}
P(\text{like} \given \text{I}) &= \frac{C(\text{I like})}{C(\text{I})} = \frac{2}{3} \\[4pt]
P(\text{hate} \given \text{I}) &= \frac{C(\text{I hate})}{C(\text{I})} = \frac{1}{3} \\[4pt]
P(\text{dogs} \given \text{like}) &= \frac{C(\text{like dogs})}{C(\text{like})} = \frac{1}{2}
\end{align*}

\subsection{The Zero-Probability Problem}

MLE has a fatal flaw: any n-gram not observed in the training corpus receives
probability \emph{zero}.  Since language is productive---speakers constantly
produce novel combinations of words---this is catastrophic.  A single unseen
n-gram in a test sentence makes the entire sentence probability zero:
\[
P(w_1, \ldots, w_N) = \prod_{t=1}^{N} P(w_t \given \ldots) = 0
\quad\text{if any factor is zero.}
\]

\begin{warningbox}
The zero-probability problem is not just a theoretical nuisance.  In speech
recognition, a zero-probability path means the recognizer will \emph{never}
produce that word sequence, regardless of how clear the acoustic evidence.
In machine translation, it means valid translations are unreachable.
Smoothing is not optional---it is essential for any n-gram model deployed in
practice.
\end{warningbox}

\subsection{Smoothing Techniques}

Smoothing methods redistribute probability mass from observed n-grams to
unobserved ones, ensuring that no event has zero probability.

\subsubsection{Laplace (Add-1) Smoothing}

The simplest approach: add 1 to every n-gram count.

\begin{mathresult}{Laplace Smoothing}
\begin{equation}
P_{\text{Laplace}}(w_t \given w_{t-1})
= \frac{C(w_{t-1}, w_t) + 1}{C(w_{t-1}) + V}
\label{eq:laplace}
\end{equation}
where $V$ is the vocabulary size.  Adding $V$ to the denominator ensures the
distribution sums to 1 over the vocabulary.
\end{mathresult}

\paragraph{Why Laplace is bad.}
For a vocabulary of $V = 100{,}000$ and a context seen $10$ times, Laplace
smoothing assigns $1/100{,}010 \approx 10^{-5}$ to each unseen bigram.
Since there are approximately $V$ unseen continuations, this steals nearly
\emph{all} the probability mass from observed events.  In practice, Laplace
smoothing wastes so much mass that it significantly degrades perplexity.

\subsubsection{Add-$k$ Smoothing}

A minor improvement: add $k < 1$ instead of $1$.
\begin{equation}
P_{\text{add-}k}(w_t \given w_{t-1})
= \frac{C(w_{t-1}, w_t) + k}{C(w_{t-1}) + kV}
\label{eq:addk}
\end{equation}
This wastes less mass but introduces a hyperparameter $k$ that must be tuned.
Still not competitive with more sophisticated methods.

\subsubsection{Interpolation and Backoff}

Rather than relying on a single n-gram order, combine multiple orders:

\paragraph{Linear interpolation.}
\begin{equation}
P_{\text{interp}}(w_t \given w_{t-2}, w_{t-1})
= \lambda_3\, P_3(w_t \given w_{t-2}, w_{t-1})
+ \lambda_2\, P_2(w_t \given w_{t-1})
+ \lambda_1\, P_1(w_t)
\label{eq:interpolation}
\end{equation}
where $\lambda_1 + \lambda_2 + \lambda_3 = 1$ and $\lambda_i \geq 0$.  The
weights $\lambda_i$ are tuned on held-out data, often via the EM algorithm.
Interpolation always has non-zero probability (because the unigram $P_1(w_t) > 0$
for all words in the vocabulary).

\paragraph{Backoff.}
Stupid backoff (Brants et al., 2007) is even simpler: use the highest-order
n-gram that has sufficient counts; if the trigram count is zero, ``back off''
to the bigram, and if that is zero, to the unigram, multiplying by a fixed
factor (0.4 in the original paper) per backoff step.  Stupid backoff does not
produce normalized probabilities---which is precisely why it scales.  The
paper's own story is the scale anchor to remember: Google trained a
\textbf{5-gram model on roughly 2 trillion tokens} of web text (about 300
billion distinct n-grams), distributed via MapReduce---far beyond what
Kneser-Ney's count-of-counts bookkeeping could handle at the time---and found
that as training data grew, the cheap unnormalized heuristic closed the
machine-translation quality gap with Kneser-Ney.  It is an early, clean
instance of a recurring lesson: enough data plus a simple method beats
sophistication on less data.  Backoff is not
inherently unnormalized, however: Katz backoff (Katz, 1987) discounts observed
n-gram counts via Good-Turing estimation and redistributes the freed mass to
the lower-order distribution, yielding a properly normalized model.

\subsubsection{Kneser-Ney Smoothing: The Gold Standard}

Kneser-Ney smoothing (Kneser \& Ney, 1995; Chen \& Goodman, 1999) is the
best-performing n-gram smoothing method.  Its key insight is that the
lower-order distribution used for backing off should not be the standard
unigram, but rather a \textbf{continuation probability}---the probability
that a word appears as a novel continuation in a new context.

\paragraph{Motivation.}
Consider the word ``Francisco.''  It is frequent in a corpus about San
Francisco, so its unigram probability $P_{\text{ML}}(\text{Francisco})$ is
high.  But Francisco almost always follows ``San''---it rarely continues
\emph{new} bigram contexts.  The continuation probability captures this:
$P_{\text{cont}}(\text{Francisco})$ is low because Francisco appears after
very few distinct words.

\begin{mathresult}{Kneser-Ney Smoothing}
\begin{align}
P_{\text{KN}}(w_t \given w_{t-1})
&= \frac{\max\bigl(C(w_{t-1}, w_t) - d,\; 0\bigr)}{C(w_{t-1})}
  + \lambda(w_{t-1})\, P_{\text{cont}}(w_t)
\label{eq:kneser_ney} \\[6pt]
P_{\text{cont}}(w_t)
&= \frac{\bigl|\{w' : C(w', w_t) > 0\}\bigr|}
        {\bigl|\{(w', w'') : C(w', w'') > 0\}\bigr|}
\label{eq:continuation}
\end{align}
where $d \in [0, 1]$ is the \textbf{absolute discount} (typically $\approx 0.75$),
and $\lambda(w_{t-1})$ is a normalizing weight that distributes the discounted
mass:
\begin{equation}
\lambda(w_{t-1}) = \frac{d}{C(w_{t-1})} \cdot \bigl|\{w : C(w_{t-1}, w) > 0\}\bigr|
\label{eq:kn_lambda}
\end{equation}
\end{mathresult}

\paragraph{Unpacking the formula.}
\begin{itemize}
    \item The first term $\frac{\max(C(w_{t-1}, w_t) - d, 0)}{C(w_{t-1})}$
          is a discounted MLE: we subtract $d$ from each observed count.
    \item The discount $d$ frees up a small amount of probability mass from
          every observed bigram.
    \item The freed mass is redistributed via $\lambda(w_{t-1}) \cdot P_{\text{cont}}(w_t)$,
          which favors words that appear in many distinct contexts (versatile words)
          over words that are frequent but context-restricted.
    \item Modified Kneser-Ney (Chen \& Goodman, 1999) uses three different discount
          values for counts of 1, 2, and $\geq 3$, yielding further improvements.
\end{itemize}

\begin{keyinsight}[Why Kneser-Ney Beats Everything Else]
The genius of Kneser-Ney is the continuation probability.  Standard backoff
uses the unigram $P(w)$ as the lower-order estimate, which rewards \emph{frequent}
words.  Kneser-Ney uses $P_{\text{cont}}(w)$, which rewards \emph{versatile}
words---those that appear after many different predecessors.  This distinction
is precisely what makes it the best n-gram smoothing method: it correctly
models the distribution of words in contexts where the higher-order n-gram
has insufficient data.
\end{keyinsight}


% ============================================================================
% SECTION 2: PERPLEXITY AND EVALUATION
% ============================================================================
\section{Perplexity and Language Model Evaluation}
\label{sec:perplexity}

How do we measure whether a language model is good?  The standard metric is
\textbf{perplexity}, which has been the primary evaluation metric for language
models from n-grams through GPT-4.

\subsection{Cross-Entropy}

Given a test corpus $W = w_1, w_2, \ldots, w_N$, the \textbf{cross-entropy}
of the model with respect to the data is:

\begin{mathresult}{Cross-Entropy of a Language Model}
\begin{equation}
H(W) = -\frac{1}{N} \sum_{t=1}^{N} \log_2 P(w_t \given w_1, \ldots, w_{t-1})
\label{eq:cross_entropy_lm}
\end{equation}
Cross-entropy measures the average number of bits needed to encode each word
under the model's probability distribution.  Lower cross-entropy means the
model assigns higher probability to the actual words that appeared.
\end{mathresult}

\subsection{Perplexity}

Perplexity is simply the exponentiation of cross-entropy:

\begin{mathresult}{Perplexity}
\begin{equation}
\text{PP}(W) = 2^{H(W)}
= P(w_1, w_2, \ldots, w_N)^{-1/N}
= \exp\!\left(-\frac{1}{N}\sum_{t=1}^{N} \ln P(w_t \given w_1, \ldots, w_{t-1})\right)
\label{eq:perplexity}
\end{equation}
where the last form uses the natural logarithm (as is standard in neural LM
evaluation).  Lower perplexity indicates a better model.

A note on base conventions: the logarithm and the exponentiation must use the
\emph{same} base.  With cross-entropy in bits (base-2 logarithms),
$\text{PP} = 2^{H}$; with cross-entropy in nats (natural logarithms),
$\text{PP} = e^{H}$.  Both conventions yield the identical perplexity value.
\end{mathresult}

\paragraph{Derivation of the equivalence.}
Starting from the definition:
\begin{align*}
\text{PP}(W) &= P(w_1, \ldots, w_N)^{-1/N} \\
&= \left(\prod_{t=1}^{N} P(w_t \given w_1, \ldots, w_{t-1})\right)^{-1/N} \\
&= \exp\!\left(-\frac{1}{N}\sum_{t=1}^{N} \ln P(w_t \given w_1, \ldots, w_{t-1})\right) \\
&= 2^{H(W)} \quad\text{(when using $\log_2$)}
\end{align*}

\subsection{Intuition: Weighted Average Branching Factor}

\begin{keyinsight}[What Perplexity Really Means]
Perplexity can be interpreted as the \textbf{weighted average number of
equally likely next words} the model considers at each position.  A perplexity
of 100 means the model is, on average, as uncertain as if it had to choose
uniformly among 100 candidates at every step.  A perfect model that always
predicts the correct next word with probability 1 has perplexity 1.  For
English text, Kneser-Ney n-gram models trained on very large corpora can reach
perplexities around 50--100, though on the small Penn Treebank benchmark a
Kneser-Ney 5-gram sits closer to 140; modern neural LMs achieve 10--30 on
standard benchmarks.  Perplexity numbers are only comparable when the training
corpus, test set, and tokenization are held fixed.
\end{keyinsight}

\paragraph{Concrete example.}
A unigram model with vocabulary $V = 10{,}000$ that assigns uniform
probability to all words has perplexity $= 10{,}000$.  A bigram model might
reduce this to $250$.  A well-trained trigram with Kneser-Ney smoothing might
achieve $80$.  GPT-2 (1.5B, evaluated zero-shot; Radford et al., 2019)
reported perplexity $35.8$ on Penn Treebank and $17.5$ on WikiText-103.

\subsection{Limitations of Perplexity}

Perplexity is useful but incomplete.  Interviewers will probe whether you
understand its limitations:

\begin{enumerate}
    \item \textbf{Vocabulary dependence.}  Perplexity is not comparable across
    models with different vocabularies.  A model with a larger vocabulary faces
    a harder prediction task, inflating perplexity.  Subword tokenization
    (BPE, WordPiece) further complicates comparison: a model that uses more
    tokens per word will have a different perplexity even on the same text.
    Normalize by computing \emph{per-character} or \emph{bits-per-byte}
    perplexity for fair comparison.

    \item \textbf{No fluency or coherence measure.}  A model can have low
    perplexity by assigning high probability to each next word \emph{locally}
    while generating globally incoherent text.  Perplexity measures
    next-token prediction, not long-range coherence.

    \item \textbf{No factuality guarantee.}  A model that confidently generates
    false statements can still have low perplexity if those statements
    resemble the training distribution.

    \item \textbf{Domain sensitivity.}  Perplexity is measured on a specific
    test set.  A model with excellent perplexity on news text may perform
    poorly on medical text.  Always report the domain.

    \item \textbf{Not a task metric.}  Low perplexity does not guarantee good
    downstream task performance.  Two models with identical perplexity can
    differ substantially on summarization, QA, or code generation.
\end{enumerate}

\begin{warningbox}
A common interview mistake: saying ``perplexity measures how good a language
model is'' without qualification.  Strong candidates immediately add: ``it
measures next-token prediction quality on a specific test set, but it
doesn't capture fluency, coherence, or factuality, and it's not comparable
across different vocabularies.''
\end{warningbox}


% ============================================================================
% SECTION 3: HIDDEN MARKOV MODELS
% ============================================================================
\section{Hidden Markov Models (HMMs)}
\label{sec:hmm}

Hidden Markov Models are \textbf{generative} probabilistic models for
sequential data where the underlying state sequence is unobserved (hidden).
For decades, HMMs were the dominant approach to part-of-speech tagging and
speech recognition (for named entity recognition---formalized only with MUC-6
in 1995---discriminative models displaced them within a few years), and they
remain the standard introduction to probabilistic sequence modeling.

\subsection{Model Specification}

An HMM is defined by:

\begin{itemize}
    \item A set of $K$ hidden states $\mathcal{S} = \{s_1, s_2, \ldots, s_K\}$.
    \item A vocabulary of $V$ possible observations $\mathcal{O} = \{o_1, o_2, \ldots, o_V\}$.
    \item \textbf{Transition probabilities} $\mA \in \R^{K \times K}$:
          $a_{ij} = P(y_t = s_j \given y_{t-1} = s_i)$.
    \item \textbf{Emission probabilities} $\mB \in \R^{K \times V}$:
          $b_j(o) = P(x_t = o \given y_t = s_j)$.
    \item \textbf{Initial state distribution} $\bm{\pi} \in \R^K$:
          $\pi_i = P(y_1 = s_i)$.
\end{itemize}

The model assumes two key \textbf{independence properties}:
\begin{enumerate}
    \item \textbf{Markov property:} The current state depends only on the
          previous state: $P(y_t \given y_1, \ldots, y_{t-1}) = P(y_t \given y_{t-1})$.
    \item \textbf{Output independence:} The observation at time $t$ depends
          only on the current state: $P(x_t \given y_1, \ldots, y_t, x_1, \ldots, x_{t-1}) = P(x_t \given y_t)$.
\end{enumerate}

\begin{figure}[H]
\centering
\begin{tikzpicture}[
    state/.style={circle, draw=deepblue, fill=lightblue, thick,
                  minimum size=1.2cm, font=\small\sffamily},
    obs/.style={rectangle, draw=deepblue!70, fill=blue!5, thick,
                minimum size=1cm, font=\small\sffamily},
    arr/.style={-{Stealth[length=5pt]}, thick, deepblue!70},
    lbl/.style={font=\scriptsize\sffamily, deepblue!70}
]
% Hidden states
\node[state] (y1) at (0,0) {$y_1$};
\node[state] (y2) at (3,0) {$y_2$};
\node[state] (y3) at (6,0) {$y_3$};
\node[state] (y4) at (9,0) {$y_4$};

% Observations
\node[obs] (x1) at (0,-2) {$x_1$};
\node[obs] (x2) at (3,-2) {$x_2$};
\node[obs] (x3) at (6,-2) {$x_3$};
\node[obs] (x4) at (9,-2) {$x_4$};

% Transitions
\draw[arr] (y1) -- node[above, lbl]{$a_{ij}$} (y2);
\draw[arr] (y2) -- node[above, lbl]{$a_{ij}$} (y3);
\draw[arr] (y3) -- node[above, lbl]{$a_{ij}$} (y4);

% Emissions
\draw[arr] (y1) -- node[right, lbl]{$b_j(o)$} (x1);
\draw[arr] (y2) -- node[right, lbl]{$b_j(o)$} (x2);
\draw[arr] (y3) -- node[right, lbl]{$b_j(o)$} (x3);
\draw[arr] (y4) -- node[right, lbl]{$b_j(o)$} (x4);

% Initial
\node[left=0.8cm of y1, lbl, font=\small\sffamily] (pi) {$\bm{\pi}$};
\draw[arr] (pi) -- (y1);

% Labels
\node[above=0.3cm of y2, font=\scriptsize\sffamily\color{gray}]
    {Hidden states (POS tags)};
\node[below=0.3cm of x2, font=\scriptsize\sffamily\color{gray}]
    {Observed words};
\end{tikzpicture}
\caption{Graphical model for an HMM.  Hidden states (top) form a first-order
Markov chain; each state emits an observation (bottom) independently.  For POS
tagging: states are tags (DET, NOUN, VERB, \ldots), observations are words.}
\label{fig:hmm_graphical}
\end{figure}

\paragraph{Joint probability of a state-observation sequence.}
Given a state sequence $\vy = (y_1, \ldots, y_T)$ and observation sequence
$\vx = (x_1, \ldots, x_T)$:
\begin{equation}
P(\vx, \vy) = \pi_{y_1} \cdot b_{y_1}(x_1) \cdot \prod_{t=2}^{T} a_{y_{t-1}, y_t} \cdot b_{y_t}(x_t)
\label{eq:hmm_joint}
\end{equation}

\subsection{The Three Fundamental HMM Problems}

Every HMM textbook organizes the theory around three problems.  Interviewers
expect you to know all three and be able to derive at least Problem 2
(Viterbi).

\subsubsection{Problem 1: Likelihood --- The Forward Algorithm}

\textbf{Given:} observation sequence $\vx = (x_1, \ldots, x_T)$ and model
parameters $(\mA, \mB, \bm{\pi})$.\\
\textbf{Find:} $P(\vx) = \sum_{\vy} P(\vx, \vy)$.

\paragraph{Naive approach.}
Enumerate all $K^T$ possible state sequences and sum.  This is
$O(K^T)$---completely intractable for even moderate $T$.

\paragraph{The forward algorithm.}
Define the \textbf{forward variable}:
\[
\alpha_t(j) = P(x_1, \ldots, x_t,\; y_t = s_j)
\]
This is the probability of observing the first $t$ observations \emph{and}
being in state $s_j$ at time $t$.

\begin{mathresult}{Forward Algorithm}
\textbf{Initialization} ($t = 1$):
\begin{equation}
\alpha_1(j) = \pi_j \cdot b_j(x_1), \qquad j = 1, \ldots, K
\label{eq:forward_init}
\end{equation}

\textbf{Recursion} ($t = 2, \ldots, T$):
\begin{equation}
\alpha_t(j) = \left[\sum_{i=1}^{K} \alpha_{t-1}(i) \cdot a_{ij}\right] \cdot b_j(x_t)
\label{eq:forward_recursion}
\end{equation}

\textbf{Termination:}
\begin{equation}
P(\vx) = \sum_{j=1}^{K} \alpha_T(j)
\label{eq:forward_termination}
\end{equation}

\textbf{Complexity:} $O(T \cdot K^2)$ time, $O(T \cdot K)$ space.
\end{mathresult}

\paragraph{Intuition.}  The recursion in Eq.~\eqref{eq:forward_recursion}
says: to compute the probability of being in state $j$ at time $t$, sum
over all states $i$ that could have preceded it, multiplying by the
transition probability $a_{ij}$ and the emission probability $b_j(x_t)$.
This is the fundamental dynamic programming insight: rather than enumerating
all paths, we accumulate probabilities incrementally.

\subsubsection{Problem 2: Decoding --- The Viterbi Algorithm}

\textbf{Given:} observation sequence $\vx$ and model parameters.\\
\textbf{Find:} the most likely state sequence
$\vy^* = \argmax_{\vy} P(\vy \given \vx)$.

The Viterbi algorithm is structurally identical to the forward algorithm,
but replaces \emph{summation} with \emph{maximization} and tracks the
argmax via backpointers.

\begin{mathresult}{Viterbi Algorithm}
Define $\delta_t(j) = \max_{y_1, \ldots, y_{t-1}} P(y_1, \ldots, y_{t-1}, y_t = s_j, x_1, \ldots, x_t)$.

\textbf{Initialization} ($t = 1$):
\begin{equation}
\delta_1(j) = \pi_j \cdot b_j(x_1), \qquad \psi_1(j) = 0
\end{equation}

\textbf{Recursion} ($t = 2, \ldots, T$):
\begin{align}
\delta_t(j) &= \max_{1 \leq i \leq K} \bigl[\delta_{t-1}(i) \cdot a_{ij}\bigr] \cdot b_j(x_t) \label{eq:viterbi_recursion} \\
\psi_t(j) &= \argmax_{1 \leq i \leq K} \bigl[\delta_{t-1}(i) \cdot a_{ij}\bigr] \label{eq:viterbi_backpointer}
\end{align}

\textbf{Termination:}
\begin{align}
y_T^* &= \argmax_{1 \leq j \leq K} \delta_T(j) \\
y_t^* &= \psi_{t+1}(y_{t+1}^*), \qquad t = T{-}1, \ldots, 1
\end{align}

\textbf{Complexity:} $O(T \cdot K^2)$ time, $O(T \cdot K)$ space.
\end{mathresult}

\paragraph{Viterbi worked example.}
Consider a tiny POS-tagging HMM with $K = 2$ states: $\{D, N\}$
(Determiner, Noun) and an observation sequence: ``the dog.''

\begin{itemize}
    \item Parameters: $\bm{\pi} = (0.6, 0.4)$, \;
          $a_{DD} = 0.3$, $a_{DN} = 0.7$, $a_{ND} = 0.4$, $a_{NN} = 0.6$.
    \item Emissions: $b_D(\text{the}) = 0.5$, $b_D(\text{dog}) = 0.1$,
          $b_N(\text{the}) = 0.1$, $b_N(\text{dog}) = 0.5$.
\end{itemize}

\textbf{Step 1: Initialization ($t=1$, word = ``the'').}
\begin{align*}
\delta_1(D) &= \pi_D \cdot b_D(\text{the}) = 0.6 \times 0.5 = 0.30 \\
\delta_1(N) &= \pi_N \cdot b_N(\text{the}) = 0.4 \times 0.1 = 0.04
\end{align*}

\textbf{Step 2: Recursion ($t=2$, word = ``dog'').}
\begin{align*}
\delta_2(D) &= \max\bigl(\delta_1(D) \cdot a_{DD},\; \delta_1(N) \cdot a_{ND}\bigr) \cdot b_D(\text{dog}) \\
            &= \max(0.30 \times 0.3,\; 0.04 \times 0.4) \times 0.1 \\
            &= \max(0.090, 0.016) \times 0.1 = 0.0090, \quad \psi_2(D) = D \\[4pt]
\delta_2(N) &= \max\bigl(\delta_1(D) \cdot a_{DN},\; \delta_1(N) \cdot a_{NN}\bigr) \cdot b_N(\text{dog}) \\
            &= \max(0.30 \times 0.7,\; 0.04 \times 0.6) \times 0.5 \\
            &= \max(0.210, 0.024) \times 0.5 = 0.1050, \quad \psi_2(N) = D
\end{align*}

\textbf{Step 3: Termination.}
$y_2^* = N$ (since $\delta_2(N) = 0.1050 > \delta_2(D) = 0.0090$),
then backtrack: $y_1^* = \psi_2(N) = D$.

\textbf{Result:} the $\to$ DET, dog $\to$ NOUN.  Exactly what we expect.

\subsubsection{The Backward Algorithm and Forward-Backward Marginals}
\label{sec:backward}

The forward variable accumulates evidence from the \emph{past}; its mirror
image, the \textbf{backward variable}, accumulates evidence from the
\emph{future}:
\[
\beta_t(i) = P(x_{t+1}, x_{t+2}, \ldots, x_T \given y_t = s_i)
\]
the probability of the remaining observations, given that we are in state
$s_i$ at time $t$.

\begin{mathresult}{Backward Algorithm}
\textbf{Initialization} ($t = T$):
\begin{equation}
\beta_T(i) = 1, \qquad i = 1, \ldots, K
\label{eq:backward_init}
\end{equation}

\textbf{Recursion} ($t = T{-}1, \ldots, 1$):
\begin{equation}
\beta_t(i) = \sum_{j=1}^{K} a_{ij} \cdot b_j(x_{t+1}) \cdot \beta_{t+1}(j)
\label{eq:backward_recursion}
\end{equation}

\textbf{Termination:}
\begin{equation}
P(\vx) = \sum_{j=1}^{K} \pi_j \cdot b_j(x_1) \cdot \beta_1(j)
\label{eq:backward_termination}
\end{equation}

\textbf{Complexity:} $O(T \cdot K^2)$ time, $O(T \cdot K)$ space---identical
to the forward pass.
\end{mathresult}

Running both recursions gives, at every position, the \textbf{posterior
marginals} that the learning algorithm (and the CRF gradient) need.  Since
$\alpha_t(i)\,\beta_t(i) = P(x_1, \ldots, x_T,\; y_t = s_i)$:

\begin{mathresult}{Forward-Backward Marginals}
\begin{align}
\gamma_t(i) &= P(y_t = s_i \given \vx)
= \frac{\alpha_t(i)\,\beta_t(i)}{P(\vx)}
\label{eq:gamma} \\[6pt]
\xi_t(i,j) &= P(y_t = s_i,\; y_{t+1} = s_j \given \vx)
= \frac{\alpha_t(i)\, a_{ij}\, b_j(x_{t+1})\, \beta_{t+1}(j)}{P(\vx)}
\label{eq:xi}
\end{align}
where $P(\vx) = \sum_i \alpha_t(i)\,\beta_t(i)$ at \emph{any} $t$---a useful
implementation sanity check.
\end{mathresult}

$\gamma_t(i)$ answers ``given the \emph{whole} sequence, how confident is the
model that position $t$ is in state $i$?''---the quantity needed for posterior
decoding and for soft counts in EM.  $\xi_t(i,j)$ is the corresponding soft
count of the transition $i \to j$ at time $t$.  The same
$\alpha \cdot \beta$ construction, run on log-space scores instead of
probabilities, yields the marginals that appear in the CRF gradient
(Section~\ref{sec:crf}).

\subsubsection{Problem 3: Learning --- The Baum-Welch Algorithm}

\textbf{Given:} observation sequences (no state labels).\\
\textbf{Find:} parameters $(\mA, \mB, \bm{\pi})$ that maximize the data
likelihood.

Baum-Welch is an instance of the \textbf{Expectation-Maximization (EM)}
algorithm, and with the forward-backward marginals in hand it is short to
state:
\begin{enumerate}
    \item \textbf{E-step:} Run forward-backward under the current parameters
    and compute the soft counts $\gamma_t(i)$ and $\xi_t(i,j)$ of
    Eqs.~\eqref{eq:gamma}--\eqref{eq:xi}.

    \item \textbf{M-step:} Re-estimate each parameter as a ratio of expected
    counts, exactly as supervised MLE would use hard counts: the new
    $\hat{a}_{ij}$ is the expected number of $i \to j$ transitions divided by
    the expected number of visits to state $i$
    ($\sum_t \xi_t(i,j) / \sum_t \gamma_t(i)$), and analogously for the
    emission and initial-state parameters.

    \item Repeat until convergence.  Each iteration provably does not decrease
    the data likelihood (the EM convergence guarantee), though the fixed point
    may be a local optimum.
\end{enumerate}

\paragraph{Practical note.}  In modern NLP, Baum-Welch is rarely used because
supervised data for POS tagging and NER is available, so parameters are
estimated directly from labeled sequences via MLE counts.  However,
understanding Baum-Welch demonstrates knowledge of EM, which is tested
independently in ML interviews.

\subsection{HMMs for POS Tagging}

The classic application: states are POS tags (DET, NOUN, VERB, ADJ, \ldots),
observations are words.

\begin{itemize}
    \item \textbf{Transition probabilities} capture syntactic patterns:
          $P(\text{NOUN} \given \text{DET})$ is high; $P(\text{DET} \given \text{DET})$
          is low.
    \item \textbf{Emission probabilities} capture lexical associations:
          $P(\text{``runs''} \given \text{VERB})$ is high;
          $P(\text{``runs''} \given \text{DET})$ is near zero.
    \item \textbf{Decoding:} Viterbi finds the globally optimal tag sequence
          $\vy^* = \argmax_{\vy} P(\vy \given \vx)$, not just the locally
          optimal tag per position.
\end{itemize}

\subsection{Limitations of HMMs}

\begin{enumerate}
    \item \textbf{Generative model assumption.}  HMMs model the joint
    $P(\vx, \vy)$, requiring strong independence assumptions about how
    observations are generated.  They cannot easily incorporate arbitrary
    overlapping features of the input.

    \item \textbf{First-order Markov assumption.}  The current tag depends
    only on the previous tag.  Richer dependencies (e.g., the tag two
    positions back, or the capitalization of the current word) are not
    naturally captured.

    \item \textbf{Output independence.}  Each word depends only on its own
    tag, not on neighboring words.  In reality, the identity of surrounding
    words is highly informative for tagging.

    \item \textbf{Discrete observations.}  Standard HMMs require discrete
    observations.  Continuous features (embeddings, character shapes) require
    extensions like Gaussian HMMs.
\end{enumerate}

These limitations directly motivate the move to Conditional Random Fields.


% ============================================================================
% SECTION 4: CONDITIONAL RANDOM FIELDS
% ============================================================================
\section{Conditional Random Fields (CRFs)}
\label{sec:crf}

Conditional Random Fields (Lafferty, McCallum, \& Pereira, 2001) are
\textbf{discriminative} structured prediction models that directly model
$P(\vy \given \vx)$ without requiring generative assumptions about the input.
CRFs address every limitation of HMMs listed above, and the CRF concept
remains actively relevant today: the \textbf{BiLSTM-CRF} (Lample et al.,
2016) and \textbf{BERT-CRF} architectures for NER use a CRF output layer
on top of neural features.

\begin{keyinsight}[Why CRFs Are a Critical Interview Topic]
The question ``Why add a CRF layer on top of BERT for NER?'' appears in a
large fraction of NLP interviews.  The answer tests whether you understand
the difference between independent per-token classification (softmax) and
structured prediction (CRF), and whether you grasp that the CRF captures
\emph{output dependencies}---constraints on valid label sequences---that the
neural encoder does not model.  This section gives you the mathematical
foundation to answer this question precisely.
\end{keyinsight}

\subsection{Generative vs.\ Discriminative: HMM vs.\ CRF}

\begin{figure}[H]
\centering
\begin{tikzpicture}[
    state/.style={circle, draw=deepblue, fill=lightblue, thick,
                  minimum size=1.1cm, font=\small\sffamily},
    obs/.style={rectangle, draw=deepblue!70, fill=blue!5, thick,
                minimum size=0.9cm, font=\small\sffamily},
    arr/.style={-{Stealth[length=5pt]}, thick, deepblue!70},
    uarr/.style={thick, deepblue!70},
    lbl/.style={font=\scriptsize\sffamily, deepblue!70}
]
% --- HMM (left) ---
\node[font=\bfseries\sffamily\color{deepblue}] at (1.5, 1.8) {HMM (Generative)};

\node[state] (hy1) at (0,0) {$y_1$};
\node[state] (hy2) at (3,0) {$y_2$};
\node[obs]   (hx1) at (0,-2) {$x_1$};
\node[obs]   (hx2) at (3,-2) {$x_2$};

\draw[arr] (hy1) -- (hy2);
\draw[arr] (hy1) -- (hx1);
\draw[arr] (hy2) -- (hx2);

\node[lbl] at (1.5, 0.4) {$P(y_t \given y_{t-1})$};
\node[lbl] at (-1.1, -1) {$P(x_t \given y_t)$};

% --- CRF (right) ---
\node[font=\bfseries\sffamily\color{deepblue}] at (8.5, 1.8) {CRF (Discriminative)};

\node[state] (cy1) at (7,0) {$y_1$};
\node[state] (cy2) at (10,0) {$y_2$};
\node[obs]   (cx1) at (7,-2) {$x_1$};
\node[obs]   (cx2) at (10,-2) {$x_2$};

\draw[uarr] (cy1) -- (cy2);
\draw[uarr] (cy1) -- (cx1);
\draw[uarr] (cy2) -- (cx2);
% CRF: x also connects to all y (global conditioning)
\draw[uarr, dashed, gray] (cx1) -- (cy2);
\draw[uarr, dashed, gray] (cx2) -- (cy1);

\node[lbl] at (8.5, 0.4) {$\psi(y_t, y_{t-1})$};
\node[lbl] at (11.2, -1) {$\psi(y_t, \vx)$};
\end{tikzpicture}
\caption{HMM vs.\ CRF graphical models.  The HMM uses directed edges (arrows)
representing generative assumptions: states generate observations.  The CRF
uses undirected edges representing compatibility functions.  Crucially, in the
CRF, each label $y_t$ can depend on the \emph{entire} input sequence $\vx$
(dashed lines), not just the local observation $x_t$.}
\label{fig:hmm_vs_crf}
\end{figure}

\begin{table}[H]
\centering
\caption{Systematic comparison of HMMs and CRFs.}
\label{tab:hmm_vs_crf}
\begin{tabularx}{\textwidth}{>{\bfseries}l X X}
\toprule
& \textbf{HMM} & \textbf{CRF} \\
\midrule
Type & Generative & Discriminative \\
Models & $P(\vx, \vy) = P(\vy) P(\vx \given \vy)$ & $P(\vy \given \vx)$ directly \\
Features & Only current observation $x_t$ & Arbitrary features of \emph{entire} input $\vx$ \\
Independence & Strong: output independence, first-order Markov & Only Markov on labels; input features are global \\
Feature overlap & Cannot use overlapping features & Naturally handles overlapping, correlated features \\
Training & MLE from counts (easy), or Baum-Welch & Gradient-based optimization (log-likelihood) \\
Decoding & Viterbi $O(TK^2)$ & Viterbi $O(TK^2)$ \\
Strengths & Simple, fast training, works with little data & Higher accuracy, flexible features, handles rich inputs \\
\bottomrule
\end{tabularx}
\end{table}

\subsection{MEMMs and the Label Bias Problem}
\label{sec:label_bias}

Between the HMM and the CRF sits a model that makes the case for CRFs sharp:
the \textbf{Maximum Entropy Markov Model} (MEMM; McCallum et al., 2000).  The
MEMM was the first widely used discriminative sequence model: it keeps the
left-to-right chain but replaces the HMM's generative factors with a
\emph{locally normalized} discriminative conditional at each step:
\begin{equation}
P(\vy \given \vx) = \prod_{t=1}^{T} P(y_t \given y_{t-1}, \vx),
\qquad
P(y_t \given y_{t-1}, \vx)
= \frac{\exp\bigl(\sum_k \lambda_k f_k(y_t, y_{t-1}, \vx, t)\bigr)}{Z_t(y_{t-1}, \vx)}
\label{eq:memm}
\end{equation}
Each factor is its own little softmax: $Z_t(y_{t-1}, \vx)$ normalizes over the
possible \emph{next} labels only.  This buys the CRF's feature flexibility---
but introduces a pathology that this chapter's own follow-up questions keep
asking about: the \textbf{label bias problem} (Lafferty et al., 2001).

\paragraph{The mechanism: local mass conservation.}
Because every per-state distribution must sum to one over its outgoing
transitions, whatever probability mass arrives at a state \emph{must} be
passed on, no matter what the observation says.  A state with few outgoing
transitions is therefore a funnel: in the extreme case of a single successor,
$P(y_{t+1} \given y_t, \vx) = 1$ \emph{regardless of the input}, and the
observation at that position is ignored outright.  Globally, paths through
low-entropy (few-successor) states are systematically favored, and per-step
conditioning on the observations cannot correct for it.

\paragraph{Worked intuition.}
Suppose the states encode which of two words we are transcribing, ``rib'' or
``rob,'' from noisy character observations (Figure~\ref{fig:label_bias}).
After the initial ``r,'' the model is on the rib branch or the rob branch, and
each branch state has exactly one successor: continue spelling its word.
Suppose training (where ``rib'' was slightly more frequent) sets the
first-step conditionals to
$P(\text{rib branch} \given s_0, \text{``r''}) = 0.55$ and
$P(\text{rob branch} \given s_0, \text{``r''}) = 0.45$.  Now decode the
observation sequence ``r,~o,~b'':
\begin{align*}
P(\text{rib path} \given \text{``r o b''})
&= 0.55 \times \underbrace{1}_{\text{on ``o''!}} \times\; 1 \;=\; 0.55 \\
P(\text{rob path} \given \text{``r o b''})
&= 0.45 \times 1 \times 1 \;=\; 0.45
\end{align*}
The middle observation is an unambiguous ``o,'' yet the MEMM decodes
``rib.''  Both branch states have a single outgoing transition, so their local
conditionals equal $1$ for \emph{any} observation---the mass each branch
carries is conserved down its funnel, and the decision was irrevocably made at
step one by the training-frequency prior.

\begin{figure}[H]
\centering
\begin{tikzpicture}[
    st/.style={circle, draw=deepblue, fill=lightblue, thick,
               minimum size=0.95cm, font=\small\sffamily},
    arr/.style={-{Stealth[length=4pt]}, thick, deepblue!70},
    lbl/.style={font=\scriptsize\sffamily, deepblue!80}
]
\node[st] (s0) at (0,0) {$s_0$};
\node[st] (r1) at (3.4,1.2) {$s_1$};
\node[st] (r2) at (3.4,-1.2) {$s_3$};
\node[st] (b1) at (6.8,1.2) {$s_2$};
\node[st] (b2) at (6.8,-1.2) {$s_4$};
\node[st] (sf) at (10.2,0) {$s_F$};

\draw[arr] (s0) -- node[lbl, above, sloped] {$0.55$ on ``r''} (r1);
\draw[arr] (s0) -- node[lbl, below, sloped] {$0.45$ on ``r''} (r2);
\draw[arr] (r1) -- node[lbl, above] {$1$ on \emph{any} $x$} (b1);
\draw[arr] (r2) -- node[lbl, below] {$1$ on \emph{any} $x$} (b2);
\draw[arr] (b1) -- node[lbl, above, sloped] {$1$ on \emph{any} $x$} (sf);
\draw[arr] (b2) -- node[lbl, below, sloped] {$1$ on \emph{any} $x$} (sf);

\node[lbl, above=0.05cm of r1, font=\scriptsize\sffamily\itshape] {``rib'' branch};
\node[lbl, below=0.05cm of r2, font=\scriptsize\sffamily\itshape] {``rob'' branch};
\end{tikzpicture}
\caption{Label bias in a MEMM.  Each branch state has a single outgoing
transition, so its locally normalized conditional is $1$ whatever the
observation; evidence seen after the branch point (the ``o'') can never change
the outcome, which is fixed at the first step.}
\label{fig:label_bias}
\end{figure}

\paragraph{The fix: global normalization.}
A CRF (next subsection) removes the per-state normalizers and replaces them
with a single partition function $Z(\vx)$ over \emph{entire label sequences}.
Potentials are unnormalized, so local mass need not be conserved: the
state--observation compatibility score at position 2 multiplies the rob path's
total score up and the rib path's down, and after the one global
normalization, the evidence seen mid-sequence can flip the decision.  All
paths compete on equal footing against all other paths, not merely against
their siblings at each state.  This---not feature flexibility, which MEMMs
already had---is the precise sense in which CRFs improve on MEMMs, and it is
the expected answer to the standard L6 probe ``what is the label bias problem
and how do CRFs avoid it?''  A modern echo: locally normalized autoregressive
models (including today's LLMs) inherit related pathologies---per-step
softmax normalization is one root of exposure bias and of degenerate
beam-search behavior---which is why globally normalized sequence models were
revisited in the neural era (Andor et al., 2016).

\subsection{Linear-Chain CRF Formulation}

A linear-chain CRF defines a conditional probability over the label sequence
$\vy = (y_1, \ldots, y_T)$ given the full input sequence
$\vx = (x_1, \ldots, x_T)$:

\begin{mathresult}{Linear-Chain CRF}
\begin{equation}
P(\vy \given \vx) = \frac{1}{Z(\vx)} \exp\!\left(\sum_{t=1}^{T} \sum_{k} \lambda_k f_k(y_t, y_{t-1}, \vx, t)\right)
\label{eq:crf}
\end{equation}
where:
\begin{itemize}
    \item $f_k(y_t, y_{t-1}, \vx, t)$ are \textbf{feature functions} that
          examine the current label, previous label, entire input, and position.
    \item $\lambda_k$ are learned weights.
    \item $Z(\vx)$ is the \textbf{partition function} that ensures normalization.
\end{itemize}
\end{mathresult}

\paragraph{Feature functions.}
Feature functions are binary or real-valued functions that capture patterns.
For a CRF used in NER, typical features include:

\begin{itemize}
    \item \textbf{Emission features:} ``current word is `London' and label is
          B-LOC'' $\to f_k = 1$.
    \item \textbf{Transition features:} ``previous label is B-PER and current
          label is I-PER'' $\to f_k = 1$.
    \item \textbf{Context features:} ``previous word is `in' and current label
          is B-LOC'' $\to f_k = 1$.
    \item \textbf{Orthographic features:} ``current word starts with uppercase
          and label is B-PER'' $\to f_k = 1$.
\end{itemize}

In the neural CRF setting (BiLSTM-CRF, BERT-CRF), the emission features are
replaced by neural network outputs, and only the \emph{transition matrix}
remains as a CRF component.

\subsection{The Partition Function}

The partition function $Z(\vx)$ sums over \emph{all possible label sequences}:

\begin{mathresult}{CRF Partition Function}
\begin{equation}
Z(\vx) = \sum_{\vy' \in \mathcal{S}^T}
\exp\!\left(\sum_{t=1}^{T} \sum_{k} \lambda_k f_k(y'_t, y'_{t-1}, \vx, t)\right)
\label{eq:partition}
\end{equation}
There are $K^T$ possible label sequences, so naive summation is intractable.
The partition function is computed efficiently using the \textbf{forward
algorithm} for CRFs, which is structurally identical to the HMM forward
algorithm but operates in log-space with scores instead of probabilities.
Complexity: $O(T \cdot K^2)$.
\end{mathresult}

\paragraph{Why the partition function matters.}
The partition function is what makes CRFs a proper probabilistic model: it
ensures $\sum_{\vy} P(\vy \given \vx) = 1$.  During training, we maximize the
log-likelihood:
\begin{equation}
\log P(\vy \given \vx) = \sum_{t=1}^{T} \sum_{k} \lambda_k f_k(y_t, y_{t-1}, \vx, t) - \log Z(\vx)
\label{eq:crf_loglik}
\end{equation}

The gradient with respect to $\lambda_k$ has the elegant form:
\begin{equation}
\frac{\partial \log P(\vy \given \vx)}{\partial \lambda_k}
= \underbrace{\sum_t f_k(y_t, y_{t-1}, \vx, t)}_{\text{observed feature count}}
- \underbrace{\E_{\vy' \sim P(\vy' | \vx)}\!\left[\sum_t f_k(y'_t, y'_{t-1}, \vx, t)\right]}_{\text{expected feature count}}
\label{eq:crf_gradient}
\end{equation}

The expected feature count requires the marginals $P(y_t, y_{t-1} \given \vx)$,
which are computed by the forward-backward algorithm for CRFs---the same
$\alpha \cdot \beta$ construction as Section~\ref{sec:backward}, run on
log-space scores instead of probabilities.

\subsection{Viterbi Decoding for CRFs}

Decoding in a CRF---finding $\vy^* = \argmax_{\vy} P(\vy \given \vx)$---uses
the Viterbi algorithm, just as in HMMs.  Since $Z(\vx)$ is constant for a
given input, the argmax depends only on the score (numerator):
\begin{equation}
\vy^* = \argmax_{\vy} \sum_{t=1}^{T} \sum_{k} \lambda_k f_k(y_t, y_{t-1}, \vx, t)
\end{equation}

The algorithm proceeds identically to the HMM Viterbi, replacing
probabilities with scores.  Complexity remains $O(T \cdot K^2)$.

\subsection{The Neural CRF: BiLSTM-CRF and BERT-CRF}
\label{sec:neural_crf}

The breakthrough of Lample et al.\ (2016) and Ma \& Hovy (2016) was combining
neural feature extractors with a CRF output layer.

\paragraph{Architecture.}
\begin{enumerate}
    \item \textbf{Feature extraction:} A BiLSTM (or BERT) processes the input
    sequence and produces a score vector $\vh_t \in \R^K$ for each position,
    where $K$ is the number of labels.  The $j$-th component $h_t^{(j)}$
    is the \emph{emission score} for label $j$ at position $t$.

    \item \textbf{Transition matrix:} A learnable matrix $\mA \in \R^{K \times K}$
    where $A_{ij}$ is the score for transitioning from label $i$ to label $j$.
    This is the CRF component.

    \item \textbf{Sequence score:}
    \begin{equation}
    s(\vx, \vy) = \sum_{t=1}^{T} h_t^{(y_t)} + \sum_{t=2}^{T} A_{y_{t-1}, y_t}
    \label{eq:neural_crf_score}
    \end{equation}

    \item \textbf{Probability:}
    \begin{equation}
    P(\vy \given \vx) = \frac{\exp\bigl(s(\vx, \vy)\bigr)}{\sum_{\vy'} \exp\bigl(s(\vx, \vy')\bigr)}
    = \frac{\exp\bigl(s(\vx, \vy)\bigr)}{Z(\vx)}
    \label{eq:neural_crf_prob}
    \end{equation}
\end{enumerate}

\paragraph{Training.}  Maximize $\log P(\vy \given \vx) = s(\vx, \vy) - \log Z(\vx)$.
The score of the gold sequence $s(\vx, \vy)$ is trivial to compute.  The
log-partition function $\log Z(\vx)$ is computed by the forward algorithm in
$O(T \cdot K^2)$.  Backpropagation flows through both the CRF layer \emph{and}
the neural encoder, enabling end-to-end training.

\paragraph{Decoding.}  Viterbi finds the best label sequence in $O(T \cdot K^2)$.

\begin{keyinsight}[Why Add CRF on Top of BERT?]
BERT captures rich \emph{input-side} dependencies: each token's representation
is informed by the entire input sequence via self-attention.  But a softmax
classifier on top of BERT makes \emph{independent} predictions for each
token---it does not model dependencies between \emph{output} labels.  The CRF
layer adds a transition matrix that models label-label constraints: for
example, it learns that I-PER should follow B-PER but not B-ORG, and that
O should not directly follow I-PER without a B-PER preceding it.  This is
especially valuable in low-data settings and for entity types with multi-token
spans.  In high-data settings with large models, the benefit is smaller
because the model learns to make consistent predictions implicitly, but the
CRF still provides a guarantee of structural validity.
\end{keyinsight}


% ============================================================================
% SECTION 5: TOPIC MODELS
% ============================================================================
\section{Topic Models: Latent Dirichlet Allocation}
\label{sec:topic_models}

\textbf{Latent Dirichlet Allocation} (LDA; Blei, Ng, \& Jordan, 2003) is
a generative probabilistic model that discovers latent ``topics'' in a
collection of documents.  While LDA has been largely superseded by neural
methods for many tasks, it remains important for interviews because:
(a) it tests understanding of generative models and Bayesian inference,
(b) it is still used in production for interpretable topic discovery and
feature engineering, and (c) modern approaches like BERTopic build on the
same conceptual framework.

\subsection{The Generative Process}

LDA posits the following generative story for a corpus of $D$ documents,
each with $N_d$ words, using $K$ topics:

\begin{enumerate}
    \item For each topic $k \in \{1, \ldots, K\}$:
    \begin{itemize}
        \item Draw a word distribution $\bm{\phi}_k \sim \text{Dir}(\bm{\beta})$
    \end{itemize}
    \item For each document $d \in \{1, \ldots, D\}$:
    \begin{itemize}
        \item Draw a topic mixture $\bm{\theta}_d \sim \text{Dir}(\bm{\alpha})$
        \item For each word position $n \in \{1, \ldots, N_d\}$:
        \begin{enumerate}
            \item Draw a topic assignment $z_{dn} \sim \text{Mult}(\bm{\theta}_d)$
            \item Draw a word $w_{dn} \sim \text{Mult}(\bm{\phi}_{z_{dn}})$
        \end{enumerate}
    \end{itemize}
\end{enumerate}

\begin{figure}[H]
\centering
\begin{tikzpicture}[
    node distance=1.5cm,
    param/.style={circle, draw=deepblue, thick,
                  minimum size=0.9cm, font=\small},
    latent/.style={circle, draw=deepblue, fill=lightblue, thick,
                   minimum size=0.9cm, font=\small},
    observed/.style={circle, draw=deepblue, fill=gray!30, thick,
                     minimum size=0.9cm, font=\small},
    hyper/.style={circle, draw=deepblue!60, thick,
                  minimum size=0.7cm, font=\small},
    plate/.style={draw=deepblue!60, thick, rounded corners=3pt,
                  inner sep=10pt},
    arr/.style={-{Stealth[length=4pt]}, thick, deepblue!70}
]
% Hyperparameters
\node[hyper] (alpha) at (0, 0) {$\bm{\alpha}$};
\node[hyper] (beta) at (8, 0) {$\bm{\beta}$};

% Document-level
\node[latent] (theta) at (2, 0) {$\bm{\theta}_d$};
\node[latent] (z) at (4, 0) {$z_{dn}$};
\node[observed] (w) at (4, -1.8) {$w_{dn}$};

% Topic-level
\node[latent] (phi) at (6, 0) {$\bm{\phi}_k$};

% Arrows
\draw[arr] (alpha) -- (theta);
\draw[arr] (theta) -- (z);
\draw[arr] (z) -- (w);
\draw[arr] (phi) -- (w);
\draw[arr] (beta) -- (phi);

% Plates
\node[plate, fit=(z)(w), label={[font=\scriptsize\sffamily, deepblue!70]below right:$N_d$}] {};
\node[plate, fit=(theta)(z)(w), label={[font=\scriptsize\sffamily, deepblue!70]below right:$D$},
      inner sep=16pt] {};
\node[plate, fit=(phi), label={[font=\scriptsize\sffamily, deepblue!70]below right:$K$}] {};

% Legend
\node[font=\scriptsize\sffamily, gray] at (4, -3.5)
    {Shaded = observed \quad Unshaded = latent \quad Plates = replication};
\end{tikzpicture}
\caption{Plate notation for Latent Dirichlet Allocation.  Each document $d$
draws a topic mixture $\bm{\theta}_d$; each word position draws a topic $z_{dn}$,
then a word $w_{dn}$ from the corresponding topic's word distribution
$\bm{\phi}_k$.  Only the words $w_{dn}$ are observed.}
\label{fig:lda_plate}
\end{figure}

\subsection{Key Quantities}

\begin{itemize}
    \item $\bm{\theta}_d \in \R^K$: document $d$'s topic mixture.
          $\theta_{dk}$ is the proportion of document $d$ devoted to topic $k$.
    \item $\bm{\phi}_k \in \R^V$: topic $k$'s word distribution.
          $\phi_{kv}$ is the probability that topic $k$ generates word $v$.
    \item $\bm{\alpha} \in \R^K$: Dirichlet prior on topic mixtures.
          Low $\alpha$ encourages documents to focus on few topics (sparse);
          high $\alpha$ encourages uniform mixtures.
    \item $\bm{\beta} \in \R^V$: Dirichlet prior on word distributions.
          Low $\beta$ encourages topics to be concentrated on few words.
\end{itemize}

\subsection{Inference}

The posterior $P(\bm{\theta}, \bm{\phi}, \mathbf{z} \given \mathbf{w}, \bm{\alpha}, \bm{\beta})$
is intractable because the partition function involves a sum over all possible
topic assignments for all words.  Two main approximation methods are used:

\paragraph{Collapsed Gibbs sampling.}
Integrate out $\bm{\theta}$ and $\bm{\phi}$ analytically (possible because the
Dirichlet is conjugate to the multinomial), then iteratively sample each
$z_{dn}$ conditioned on all other assignments:
\begin{equation}
P(z_{dn} = k \given \mathbf{z}_{-dn}, \mathbf{w})
\propto \frac{n_{dk}^{-dn} + \alpha_k}{\sum_{k'} (n_{dk'}^{-dn} + \alpha_{k'})}
\cdot \frac{n_{kv}^{-dn} + \beta_v}{\sum_{v'} (n_{kv'}^{-dn} + \beta_{v'})}
\label{eq:gibbs_lda}
\end{equation}
where $n_{dk}^{-dn}$ is the count of words in document $d$ assigned to topic
$k$ (excluding position $n$), and $n_{kv}^{-dn}$ is the count of word $v$
assigned to topic $k$ (excluding position $n$).

\paragraph{Variational inference.}
Approximate the posterior with a factored variational distribution and
optimize the Evidence Lower Bound (ELBO).  Faster than Gibbs for large
corpora and amenable to online/stochastic updates.

\subsection{Modern Relevance: BERTopic}

Classical LDA represents documents as bags of words, losing all semantic and
contextual information.  \textbf{BERTopic} (Grootendorst, 2022) modernizes
topic modeling:

\begin{enumerate}
    \item Embed documents using a pretrained sentence encoder (e.g., all-MiniLM).
    \item Reduce dimensionality with UMAP.
    \item Cluster with HDBSCAN (density-based, discovers variable numbers of topics).
    \item Extract topic representations using class-based TF-IDF (c-TF-IDF).
\end{enumerate}

BERTopic produces more coherent topics than LDA on most benchmarks, handles
short texts better (no bag-of-words sparsity), and does not require specifying
$K$ in advance.  However, LDA remains preferred when interpretable generative
assumptions are needed or when Bayesian posterior inference over topics is
required.


% ============================================================================
% SECTION 6: STATISTICAL MACHINE TRANSLATION
% ============================================================================
\section{Statistical Machine Translation}
\label{sec:smt}

Before neural machine translation (NMT), the dominant paradigm was
\textbf{statistical machine translation (SMT)}.  Understanding SMT is
valuable not because it is used in production today, but because:
(a) it introduces the noisy channel model, a foundational concept in
probabilistic NLP; (b) it contextualizes why the shift to neural MT was
revolutionary; and (c) interviewers at companies with translation products
(Google, Meta, Apple, Amazon) sometimes ask about the history.

\subsection{The Noisy Channel Model}

The central idea: treat translation as a decoding problem.  Given a source
sentence $\mathbf{f}$ (foreign), find the target sentence $\mathbf{e}$
(English) that maximizes:

\begin{mathresult}{Noisy Channel for Translation}
\begin{equation}
\mathbf{e}^* = \argmax_{\mathbf{e}}\; P(\mathbf{e} \given \mathbf{f})
= \argmax_{\mathbf{e}}\; \underbrace{P(\mathbf{f} \given \mathbf{e})}_{\text{translation model}} \cdot \underbrace{P(\mathbf{e})}_{\text{language model}}
\label{eq:noisy_channel}
\end{equation}
The equality follows from Bayes' theorem; $P(\mathbf{f})$ is constant and
drops from the argmax.
\end{mathresult}

\paragraph{Why decompose this way?}
\begin{itemize}
    \item The \textbf{language model} $P(\mathbf{e})$ ensures fluency---it
    prefers grammatically correct English, and can be trained on large
    monolingual corpora.
    \item The \textbf{translation model} $P(\mathbf{f} \given \mathbf{e})$
    ensures adequacy---it models word-level and phrase-level translation
    equivalences.  It is trained on parallel (bilingual) corpora.
    \item Separating the two allows each to be improved independently.
\end{itemize}

\subsection{Alignment and IBM Models}

The translation model requires modeling \textbf{word alignment}: which source
words correspond to which target words.  The IBM Models (Brown et al., 1993)
define a series of increasingly complex alignment models:

\begin{itemize}
    \item \textbf{IBM Model 1:} Uniform alignment probabilities.  All alignments
    equally likely.  Only learns word-level translation probabilities $t(f \given e)$.
    Tractable via EM.
    \item \textbf{IBM Model 2:} Adds position-dependent alignment probabilities
    $a(i \given j, l_e, l_f)$: words near the diagonal of the sentence pair
    align more often.
    \item \textbf{IBM Models 3--5:} Historical one-liners today: they added
    fertility (one English word generating multiple foreign words) and
    distortion modeling at steeply increasing implementation complexity.  Know
    that they exist and what fertility means; 2026 interviews do not go deeper.
\end{itemize}

\paragraph{Phrase-based SMT.}
The practical workhorse (Koehn et al., 2003) extended word-level alignment
to \textbf{phrase pairs}: contiguous sequences of words that translate
together.  A phrase table stores translation probabilities for millions of
phrase pairs, and a log-linear model combines the phrase translation
probability, language model score, reordering model, and word penalty.
Beam search decodes the best translation.

\subsection{Why Neural MT Replaced SMT}

The transition from SMT to NMT (Sutskever et al., 2014; Bahdanau et al., 2015)
was driven by several factors:

\begin{enumerate}
    \item \textbf{End-to-end training.}  NMT trains a single model end-to-end,
    while SMT requires separately training a translation model, language model,
    reordering model, and then combining them with a log-linear model.

    \item \textbf{Fluency.}  NMT generates more fluent output because it
    models the full target-side context, not just n-gram language model
    scores over independently selected phrases.

    \item \textbf{Handling of long-range dependencies.}  SMT reordering models
    are local; NMT attention mechanisms can capture long-distance dependencies.

    \item \textbf{Morphological richness.}  SMT struggles with morphologically
    complex languages because phrase tables become enormous.  NMT with subword
    tokenization handles this naturally.

    \item \textbf{Simplicity.}  A single encoder-decoder model replaces an
    elaborate pipeline of components.
\end{enumerate}

\begin{warningbox}
Do not dismiss SMT entirely in an interview.  If asked ``why did NMT replace
SMT?'' a strong candidate explains the specific advantages (end-to-end,
fluency, long-range dependencies) rather than just saying ``neural is better.''
Showing respect for the engineering of SMT systems---which served billions of
translations at Google, Microsoft, and others for over a decade---signals
intellectual maturity.
\end{warningbox}


% ============================================================================
% SECTION 7: MODERN PAYOFFS
% ============================================================================
\section{Modern Payoffs of Statistical NLP}
\label{sec:modern_payoffs}

Why should a 2026 candidate spend preparation hours on n-gram counts, lattices,
and dynamic programming over label sequences?  Because the machinery did not
retire with the models---it re-emerged inside the LLM stack.  This section is
the transfer map.

\subsection{Constrained Decoding: Automata Meet LLMs}

Structured-output enforcement---making an LLM emit valid JSON, conform to a
schema, or follow a grammar---is this chapter's lattice machinery reborn.
Conceptually, the constraint (a regular expression, JSON schema, or
context-free grammar) is compiled into a finite-state automaton, and
generation is the \emph{intersection} of that automaton with the LLM's
next-token distribution: at each step the decoder masks every token that would
leave the automaton's set of reachable accepting paths, renormalizes, and
samples.  Viewed as weighted-automaton (WFSA) intersection, this is the same
composition of scoring machines that powered SMT and speech decoders for two
decades---the LLM supplies one set of weights, the constraint automaton
supplies hard zeroes.  The serving-side implementation (aligning automaton
states with a subword vocabulary, batched token masking) belongs to the Deep
Learning volume's inference-optimization chapter [DL 19].

The dynamic-programming skills transfer the same way: beam search over an
LLM's token lattice is approximate Viterbi with a width-$k$ frontier, and CTC
training in speech recognition is the forward-backward algorithm of
Section~\ref{sec:backward} applied to an alignment lattice.  When an
interviewer watches you derive Viterbi, this transfer is what they are
actually pricing.

\subsection{Infini-gram: N-gram Models over Trillions of Tokens}

The counting machinery of Section~\ref{sec:ngram_lm} has been rebuilt at
modern scale.  \textbf{Infini-gram} (Liu et al., 2024) drops precomputed count
tables entirely: it indexes a trillion-token corpus with a suffix array, so
the count of \emph{any} n-gram---for unbounded $n$---is retrieved in
milliseconds at query time.  The resulting ``$\infty$-gram'' language model
backs off to the \emph{longest} context suffix present in the corpus rather
than to a fixed order.  Its 2026 role is not generation but \emph{analysis and
attribution}: quantifying how much of an LLM's output is verbatim corpus
n-grams, tracing generated text back to candidate training documents,
auditing benchmark contamination and membership, and---interpolated with a
neural LM---localizing exactly where neural probabilities diverge from corpus
statistics.

\subsection{Statistical LMs as Data Instruments}

The most consequential production role of the n-gram LM today is
\emph{upstream} of the neural model: KenLM perplexity filters score every
candidate pretraining document against a clean-corpus 5-gram model, while
MinHash-over-shingles deduplication and 13-gram decontamination guard the
corpus itself.  The full treatment---including the fastText quality-classifier
pattern and a worked 10T-token pipeline design---is in
Section~\ref{sec:classical_llm_stack} of
Chapter~\ref{chap:classical_foundations}; we do not repeat it here.  One
final one-liner completes the map: cheap n-gram models also serve as
\emph{draft models for speculative decoding}, proposing continuation tokens
that the large LLM verifies in parallel, because their per-token cost is
effectively zero.


% ============================================================================
% SECTION 8: INTERVIEW QUESTIONS
% ============================================================================
\section{Interview Questions}
\label{sec:stat_nlp_questions}

\begin{interviewq}
\textbf{Q: Explain the Viterbi algorithm.  What is its time complexity?
Trace through a small example.}

\badgeLfive{L5} \quad \badgetype{Mathematical} \quad \badgefreq{\freqmed}

\stronganswer
The Viterbi algorithm finds the most likely sequence of hidden states given a
sequence of observations in a probabilistic sequence model (HMM or CRF).  It
is a dynamic programming algorithm that replaces the naive enumeration of all
$K^T$ state sequences with an efficient recursion.

The key idea is to define $\delta_t(j)$ as the probability of the most likely
path ending in state $j$ at time $t$.  The recursion is:
$\delta_t(j) = \max_i [\delta_{t-1}(i) \cdot a_{ij}] \cdot b_j(x_t)$, where
$a_{ij}$ is the transition probability and $b_j(x_t)$ is the emission probability.
We also store backpointers $\psi_t(j) = \argmax_i [\delta_{t-1}(i) \cdot a_{ij}]$
to recover the optimal path by tracing back from the best final state.

The time complexity is $O(T \cdot K^2)$: at each of $T$ positions, we consider
$K$ current states, and for each we maximize over $K$ predecessor states.
Space is $O(T \cdot K)$ for the backpointer table.

\emph{[Candidate should then trace through a 2-state, 2-observation example
as shown in Section~\ref{sec:hmm}, computing $\delta$ values and backpointers
step by step.]}

\redflags
\begin{itemize}
    \item Confusing Viterbi (max, gives best path) with the forward algorithm
          (sum, gives total probability).
    \item Stating the complexity as $O(T \cdot K)$ instead of $O(T \cdot K^2)$.
    \item Inability to trace through a concrete example.
    \item Forgetting the backpointer step---just computing $\delta_T$ without
          explaining how to recover the path.
\end{itemize}

\interviewtest{Whether the candidate understands dynamic programming, can
distinguish between computing the best path and the total probability, and can
apply the algorithm concretely.}

\goodgreat
\begin{itemize}
    \item \badgeLfive{L5}: Correctly states the algorithm, complexity, and traces
          through an example.  Mentions that it applies to both HMMs and CRFs.
    \item \badgeLsix{L6}: Discusses numerical issues (working in log-space to
          avoid underflow), relationship to the forward algorithm, and how
          Viterbi is used in the BiLSTM-CRF decoding pipeline.
    \item \badgeLseven{L7}: Connects to beam search as an approximate
          alternative for models where exact Viterbi is intractable, discusses
          the Viterbi algorithm for higher-order CRFs ($O(T \cdot K^{n+1})$
          for $n$-th order), and mentions A* search variants.
\end{itemize}

\followups{What is the forward algorithm and how does it differ from Viterbi?
How do you handle numerical underflow?  How is Viterbi used in CRFs?
What happens to the complexity for a second-order model?}
\end{interviewq}

\vspace{0.5em}

\begin{interviewq}
\textbf{Q: Why use a CRF layer on top of BERT for NER instead of just a
softmax classifier?}

\badgeLfive{L5} \quad \badgetype{Conceptual} \quad \badgefreq{\freqhigh}

\stronganswer
A softmax classifier on top of BERT makes \emph{independent} predictions for
each token.  At position $t$, it computes
$P(y_t \given \vx) = \softmax(\mW \vh_t + \vb)$, where $\vh_t$ is BERT's
contextualized representation.  This means the prediction at position $t$
does not consider what labels were assigned to neighboring positions.

A CRF layer adds a transition score matrix $\mA \in \R^{K \times K}$ that
models the compatibility between adjacent labels.  The probability of the
entire sequence becomes
$P(\vy \given \vx) \propto \exp(\sum_t h_t^{(y_t)} + \sum_t A_{y_{t-1}, y_t})$.

This captures output-side dependencies: the CRF learns that I-PER should
follow B-PER but not B-ORG, that O cannot directly follow an I- tag without a
preceding B- tag, and so on.  Without the CRF, a softmax classifier might
produce invalid label sequences like ``O I-PER I-ORG,'' which the CRF
transition matrix would prevent.

The benefit is most pronounced in: (a) low-data regimes where the model has
not seen enough examples to learn consistent labeling patterns implicitly;
(b) tasks with many entity types and complex tag schemes (BIOES); (c) when
the model is smaller or less capable.

\redflags
\begin{itemize}
    \item Saying ``BERT already captures dependencies, so CRF is unnecessary''
          without distinguishing input-side from output-side dependencies.
    \item Not knowing what a CRF transition matrix does.
    \item Being unable to explain what ``invalid label sequences'' means.
    \item Saying the CRF always helps---in practice, for large fine-tuned BERT
          models with abundant data, the improvement is often marginal.
\end{itemize}

\interviewtest{Whether the candidate understands structured prediction,
distinguishes input-side from output-side modeling, and can reason about when
additional model complexity is justified.}

\goodgreat
\begin{itemize}
    \item \badgeLfive{L5}: Clearly explains the independent-vs-structured
          distinction, gives examples of invalid sequences.
    \item \badgeLsix{L6}: Explains the training procedure (forward algorithm
          for $\log Z$, Viterbi for decoding), discusses when CRF helps vs.\
          when softmax suffices, and mentions computational overhead ($O(TK^2)$
          per training example for the CRF forward pass).
    \item \badgeLseven{L7}: Discusses span-based NER as an alternative to
          sequence labeling entirely, nested NER, the label bias problem in
          MEMMs that CRFs avoid, and production tradeoffs (CRF adds latency
          during inference via Viterbi).
\end{itemize}

\followups{How is the CRF layer trained end-to-end with BERT?  What is the
partition function and how is it computed?  When does CRF not help?
What is the label bias problem?}
\end{interviewq}

\vspace{0.5em}

\begin{interviewq}
\textbf{Q: What is perplexity?  How does it relate to cross-entropy?
What are its limitations as an evaluation metric?}

\badgeLfive{L5} \quad \badgetype{Mathematical} \quad \badgefreq{\freqhigh}

\stronganswer
Perplexity measures how well a probabilistic language model predicts a test
set.  It is defined as:
$\text{PP}(W) = \exp\!\bigl(-\frac{1}{N}\sum_{t=1}^{N} \ln P(w_t \given w_1, \ldots, w_{t-1})\bigr)$,
which equals $2^{H(W)}$ where $H(W)$ is the cross-entropy per word.

Cross-entropy $H(W) = -\frac{1}{N}\sum_t \log_2 P(w_t \given \text{context})$
measures the average number of bits needed to encode each word under the model.
Perplexity is the exponentiation of this quantity, giving it an intuitive
interpretation as the weighted average number of equally likely next words the
model considers at each step.  Lower perplexity means the model is more certain
and assigns higher probability to the actual words.

Limitations: (1) Perplexity is vocabulary-dependent---models with different
tokenizations are not directly comparable (use bits-per-byte for fair
comparison).  (2) It only measures local next-token prediction, not global
coherence or factual correctness.  (3) It is domain-specific---a model can
have excellent perplexity on news but poor perplexity on medical text.  (4) It
does not predict downstream task performance: two models with similar perplexity
can differ substantially on QA or summarization.

\redflags
\begin{itemize}
    \item Not knowing the formula or the connection to cross-entropy.
    \item Saying ``lower perplexity is always better'' without qualification.
    \item Claiming perplexity measures fluency, coherence, or factuality.
    \item Not knowing that perplexity is not comparable across different
          vocabularies or tokenizers.
\end{itemize}

\interviewtest{Mathematical precision, understanding of information-theoretic
foundations, and awareness that metrics have limitations.}

\goodgreat
\begin{itemize}
    \item \badgeLfive{L5}: States the formula, gives the branching factor
          interpretation, lists 2--3 limitations.
    \item \badgeLsix{L6}: Derives the relationship between cross-entropy and
          perplexity, discusses bits-per-byte for fair comparison, connects to
          the KL divergence between the true distribution and the model.
    \item \badgeLseven{L7}: Discusses perplexity in the context of LLM
          evaluation (why labs report perplexity alongside benchmark scores),
          the gap between perplexity and human preference rankings, and
          alternative metrics like calibration and task-specific evaluations.
\end{itemize}

\followups{What is the perplexity of a uniform model over vocabulary $V$?
How do you compute perplexity for subword models?  Why does GPT-4's
perplexity not tell us how good it is at reasoning?}
\end{interviewq}

\vspace{0.5em}

\begin{interviewq}
\textbf{Q: Compare generative and discriminative sequence models.  Why are
CRFs generally preferred over HMMs for sequence labeling?}

\badgeLfive{L5} \quad \badgetype{Conceptual} \quad \badgefreq{\freqmed}

\stronganswer
An HMM is a \emph{generative} model: it models the joint probability
$P(\vx, \vy) = P(\vy) \cdot P(\vx \given \vy)$ by specifying how hidden states
generate observations.  This requires strong independence assumptions: each
observation depends only on its own hidden state (output independence), and
the transition structure is first-order Markov.

A CRF is a \emph{discriminative} model: it directly models the conditional
$P(\vy \given \vx)$ without making assumptions about how the inputs were
generated.  This has three key advantages:

First, CRFs can use \emph{arbitrary, overlapping features} of the entire input.
In NER, useful features include the current word, its capitalization, the
previous and next words, prefix/suffix patterns, and gazetteer membership.  An
HMM cannot naturally incorporate such features because its emission model is
$P(\text{word} \given \text{tag})$---a single discrete distribution.

Second, CRFs avoid the \emph{label bias problem}
(Section~\ref{sec:label_bias}) that affects Maximum Entropy Markov Models
(MEMMs), which are locally normalized.  CRFs are globally normalized via the
partition function $Z(\vx)$, ensuring that all paths compete on equal footing.

Third, CRFs consistently achieve higher accuracy on sequence labeling
benchmarks because they model $P(\vy \given \vx)$ directly, optimizing what
we actually care about, rather than allocating modeling capacity to the input
distribution $P(\vx)$.

CRFs are preferred in practice, but HMMs retain advantages in: (a) unsupervised
settings (Baum-Welch enables learning without labels), (b) very low-data
scenarios (generative models can be more data-efficient), and (c) when you
need to generate observations (HMMs are generative; CRFs are not).

\redflags
\begin{itemize}
    \item Not knowing the distinction between generative and discriminative.
    \item Saying ``HMMs are obsolete''---they are still the foundation of
          understanding sequence models.
    \item Not mentioning the feature flexibility advantage of CRFs.
    \item Confusing MEMMs with CRFs.
\end{itemize}

\interviewtest{Whether the candidate understands the generative/discriminative
distinction at a fundamental level and can reason about when each paradigm
is appropriate.}

\goodgreat
\begin{itemize}
    \item \badgeLfive{L5}: Clearly states the distinction, gives the feature
          flexibility argument, and mentions CRFs' higher accuracy.
    \item \badgeLsix{L6}: Explains the label bias problem, discusses global
          vs.\ local normalization, and connects to the Ng \& Jordan (2001)
          result on generative vs.\ discriminative classifiers.
    \item \badgeLseven{L7}: Discusses how the same generative/discriminative
          dichotomy appears in modern NLP (e.g., masked language models as
          generative feature extractors, discriminative fine-tuning), and
          when hybrid approaches work.
\end{itemize}

\followups{What is the label bias problem?  Why does global normalization help?
When would you prefer a generative model?  How does this distinction relate to
BERT vs.\ GPT?}
\end{interviewq}

\vspace{0.5em}

\begin{interviewq}
\textbf{Q: Explain Kneser-Ney smoothing.  Why is it better than add-1
(Laplace) smoothing?}

\badgeLsix{L6} \quad \badgetype{Mathematical} \quad \badgefreq{\freqlow}

\stronganswer
The core problem with n-gram language models is that any n-gram not seen in
training gets zero probability under MLE.  Smoothing redistributes probability
mass from observed to unobserved events.

Laplace (add-1) smoothing adds 1 to every count:
$P(w \given h) = (C(h, w) + 1) / (C(h) + V)$.  The problem is that for
large vocabularies ($V \approx 100{,}000$), this steals enormous probability
mass from observed events and distributes it uniformly.  It treats all unseen
n-grams as equally likely, which is wrong.

Kneser-Ney smoothing improves on this with two innovations.  First,
\emph{absolute discounting}: instead of adding pseudocounts, it subtracts a
fixed discount $d \approx 0.75$ from each observed count, freeing a controlled
amount of mass.  Second, and more importantly, the \emph{continuation
probability}: the lower-order distribution used for backoff is not the standard
unigram $P(w) = C(w)/N$ but
$P_{\text{cont}}(w) = |\{v : C(v, w) > 0\}| / |\{(v, v') : C(v, v') > 0\}|$,
which measures how many distinct contexts the word continues.

The classic example: ``Francisco'' is frequent but almost always follows
``San.''  Standard unigram would give it high probability as a backoff.
But $P_{\text{cont}}(\text{Francisco})$ is low because it appears after very
few distinct predecessors.  Kneser-Ney correctly assigns low probability to
``Francisco'' appearing after a novel context.

Modified Kneser-Ney (Chen \& Goodman, 1999) uses three different discount
values for counts of 1, 2, and $\geq 3$, and is the best-performing n-gram
smoothing method known.

\redflags
\begin{itemize}
    \item Not knowing what continuation probability is.
    \item Confusing Kneser-Ney with simple interpolation or backoff.
    \item Saying ``Kneser-Ney just subtracts a discount'' without explaining
          the continuation probability, which is the key innovation.
    \item Not being able to give an example like ``Francisco.''
\end{itemize}

\interviewtest{Depth of statistical NLP knowledge, ability to explain why
a sophisticated method works better than a simple one, and whether the
candidate has studied classical NLP beyond surface level.}

\goodgreat
\begin{itemize}
    \item \badgeLfive{L5}: Explains absolute discounting and continuation
          probability with an example.
    \item \badgeLsix{L6}: Writes the formulas for both terms, explains how
          the discount $d$ is estimated, and discusses modified Kneser-Ney.
    \item \badgeLseven{L7}: Connects to the information-theoretic perspective
          (smoothing as prior), discusses how Kneser-Ney relates to Bayesian
          n-gram models (Pitman-Yor processes), and notes that neural LMs
          implicitly smooth via continuous representations.
\end{itemize}

\followups{What is absolute discounting?  How is the discount $d$ typically
chosen?  What is modified Kneser-Ney?  How do neural language models handle
the zero-probability problem differently?}
\end{interviewq}

\vspace{0.5em}

\begin{interviewq}
\textbf{Q: Explain LDA's generative process.  What are the latent variables
and how is inference performed?}

\badgeLfive{L5} \quad \badgetype{Conceptual} \quad \badgefreq{\freqlow}

\stronganswer
LDA (Latent Dirichlet Allocation) is a generative probabilistic model for
discovering topics in a document corpus.  The generative process is:

For each topic $k$, draw a word distribution $\bm{\phi}_k$ from a Dirichlet
prior $\text{Dir}(\bm{\beta})$.  For each document $d$, draw a topic mixture
$\bm{\theta}_d$ from $\text{Dir}(\bm{\alpha})$.  Then, for each word position
in the document, first sample a topic $z$ from the multinomial
$\text{Mult}(\bm{\theta}_d)$, then sample a word $w$ from
$\text{Mult}(\bm{\phi}_z)$.

The latent variables are: (1) the per-document topic mixtures $\bm{\theta}_d$,
(2) the per-topic word distributions $\bm{\phi}_k$, and (3) the per-word topic
assignments $z_{dn}$.  Only the words $w_{dn}$ are observed.

The posterior is intractable because the partition function requires summing
over all possible topic assignments for all words.  Two main inference methods
are: (a) collapsed Gibbs sampling---integrate out $\bm{\theta}$ and
$\bm{\phi}$ analytically, then iteratively sample each $z_{dn}$; and
(b) variational inference---approximate the posterior with a factored
distribution and optimize the ELBO.

The hyperparameter $\alpha$ controls topic sparsity (low $\alpha$ means
documents focus on few topics), and $\beta$ controls word sparsity within
topics (low $\beta$ means topics are concentrated on few words).

\redflags
\begin{itemize}
    \item Confusing the generative direction (topics generate words, not the
          other way around).
    \item Not knowing the role of the Dirichlet distribution.
    \item Saying ``LDA uses backpropagation'' (it does not---it predates
          deep learning).
    \item Not distinguishing observed from latent variables.
\end{itemize}

\interviewtest{Understanding of generative models, Bayesian reasoning,
and approximate inference.  Tests whether the candidate can think beyond
discriminative neural models.}

\goodgreat
\begin{itemize}
    \item \badgeLfive{L5}: Correctly describes the generative process and
          identifies all latent variables.
    \item \badgeLsix{L6}: Explains why the posterior is intractable, describes
          Gibbs sampling or variational inference concretely, and discusses
          hyperparameter effects.
    \item \badgeLseven{L7}: Connects to modern topic modeling (BERTopic),
          discusses limitations of the bag-of-words assumption, and mentions
          extensions like the correlated topic model or dynamic topic models.
\end{itemize}

\followups{Why is the Dirichlet distribution used as the prior?  What happens
if $\alpha$ is very large?  How does BERTopic differ from LDA?  When would you
use LDA today?}
\end{interviewq}

\vspace{0.5em}

\begin{interviewq}
\textbf{Q: How do n-gram language models relate to modern neural language
models?  What did each generation preserve and what did it change?}

\badgeLfive{L5} \quad \badgetype{Conceptual} \quad \badgefreq{\freqmed}

\stronganswer
Both n-gram models and neural language models solve the same fundamental
problem: estimating $P(w_t \given w_1, \ldots, w_{t-1})$, the probability of
the next word given the context.  The chain rule decomposition is identical.
What changes is how the conditional is parameterized.

N-gram models use the Markov assumption to truncate context to $n{-}1$ words,
then estimate probabilities via smoothed count ratios.  Neural LMs (Bengio
et al., 2003) replaced the count table with a neural network that maps the
context to a probability distribution over the vocabulary.

The evolution preserves and extends:
\begin{itemize}
    \item \emph{Preserved:} The chain rule decomposition; the left-to-right
    autoregressive factorization; perplexity as the evaluation metric;
    the fundamental goal of next-word prediction.
    \item \emph{Changed:} Fixed context window $\to$ unbounded context (RNNs,
    then transformers); discrete count tables $\to$ continuous parameter
    spaces; manual smoothing $\to$ implicit smoothing through continuous
    representations; sparse features $\to$ dense embeddings.
\end{itemize}

The key insight is that neural LMs solved the \emph{curse of dimensionality}
that plagues n-grams: for large $n$, most n-grams are never observed.
By mapping words to continuous embeddings, neural LMs share statistical
strength between similar words---seeing ``the cat sat on the mat'' helps
predict ``the dog sat on the rug'' because ``cat'' and ``dog'' have similar
embeddings.  N-gram models treat these as completely independent events.

Modern LLMs (GPT, LLaMA, Claude) are still autoregressive language models,
just with transformer architectures, massive scale, and alignment training.
The intellectual lineage from n-gram LMs is direct and unbroken.

\redflags
\begin{itemize}
    \item Saying ``n-gram models are completely different from neural LMs''
          without seeing the shared foundation.
    \item Not knowing the chain rule decomposition.
    \item Being unable to articulate the curse of dimensionality that neural
          models solve.
    \item Treating the history as irrelevant rather than as a progression.
\end{itemize}

\interviewtest{Whether the candidate sees NLP as a coherent field with an
intellectual arc, not just a collection of disconnected techniques.  Tests
depth of understanding of language modeling as a concept.}

\goodgreat
\begin{itemize}
    \item \badgeLfive{L5}: Identifies the shared chain rule foundation and
          the key differences (continuous representations, longer context).
    \item \badgeLsix{L6}: Discusses Bengio (2003) as the bridge, explains the
          curse of dimensionality, and connects smoothing to regularization in
          neural models.
    \item \badgeLseven{L7}: Discusses scaling laws showing that both n-gram
          and neural LMs follow power-law improvements with data, connects
          to the entropy rate of language as a fundamental limit, and discusses
          how modern LLMs are still evaluated by perplexity.
\end{itemize}

\followups{What was the first neural language model?  Why did it take so long
for neural LMs to dominate?  What is the entropy rate of English?
Do modern LLMs still use the Markov assumption?}
\end{interviewq}

\vspace{0.5em}

\begin{interviewq}
\textbf{Q: What is the partition function in a CRF?  Why is it
computationally challenging, and how is it computed efficiently?}

\badgeLsix{L6} \quad \badgetype{Mathematical} \quad \badgefreq{\freqmed}

\stronganswer
The partition function $Z(\vx) = \sum_{\vy'} \exp(s(\vx, \vy'))$ normalizes
the CRF distribution, ensuring $\sum_{\vy} P(\vy \given \vx) = 1$.  The sum
is over all possible label sequences $\vy'$.

The computational challenge is that there are $K^T$ possible sequences for $K$
labels and $T$ positions.  For a typical NER task with $K = 17$ labels (BIOES
for 4 entity types plus O) and $T = 100$ tokens, that is $17^{100}$
sequences---an astronomically large number.  Naive enumeration is impossible.

The solution is the \textbf{forward algorithm} for CRFs, which exploits the
linear-chain structure.  Define forward scores
$\alpha_t(j) = \log \sum_{\vy_{1:t}: y_t=j} \exp(s(\vx, \vy_{1:t}))$.
The recursion is:
$\alpha_t(j) = \text{logsumexp}_i[\alpha_{t-1}(i) + A_{ij}] + h_t^{(j)}$,
where $A_{ij}$ is the transition score and $h_t^{(j)}$ is the emission score.
Then $\log Z(\vx) = \text{logsumexp}_j[\alpha_T(j)]$.

This runs in $O(T \cdot K^2)$ time, making the partition function tractable
for any practical sequence length and label set.  The same algorithm is used
during training (to compute $\log Z$ for the loss) and the forward-backward
algorithm extends it to compute the marginals needed for the gradient.

\redflags
\begin{itemize}
    \item Not knowing what the partition function is or why it is needed.
    \item Saying the complexity is $O(K^T)$ without knowing the dynamic
          programming solution.
    \item Confusing the forward algorithm (sum, for $Z$) with Viterbi
          (max, for decoding).
    \item Not mentioning that computations are done in log-space to avoid
          numerical underflow.
\end{itemize}

\interviewtest{Mathematical maturity, understanding of structured prediction
and dynamic programming, and familiarity with the training mechanics of
CRF-based models.}

\goodgreat
\begin{itemize}
    \item \badgeLfive{L5}: States the definition, explains why naive computation
          is intractable, and describes the forward algorithm at a high level.
    \item \badgeLsix{L6}: Writes the forward recursion explicitly, explains
          log-space computation (logsumexp trick), and connects to the gradient
          computation (observed minus expected feature counts).
    \item \badgeLseven{L7}: Discusses generalizations---partition functions in
          higher-order CRFs, semi-Markov CRFs (for segment-level models),
          approximate methods when exact computation is intractable (e.g.,
          belief propagation in general CRFs), and connections to the
          partition function in statistical physics.
\end{itemize}

\followups{How does the gradient of the CRF log-likelihood relate to the
partition function?  What is the forward-backward algorithm?
What happens to the complexity for a second-order CRF?  How does this
relate to the partition function in energy-based models?}
\end{interviewq}

\vspace{0.5em}

\begin{interviewq}
\textbf{Q: Explain the noisy channel model for machine translation.  Why
was it eventually replaced by neural MT?}

\badgeLfive{L5} \quad \badgetype{Conceptual} \quad \badgefreq{\freqlow}

\stronganswer
The noisy channel model frames translation as Bayesian inference.  To
translate a foreign sentence $\mathbf{f}$ to English, we find:
$\mathbf{e}^* = \argmax_{\mathbf{e}} P(\mathbf{f} \given \mathbf{e}) \cdot P(\mathbf{e})$,
via Bayes' theorem.

This decomposes into two independently trainable components: a
\emph{translation model} $P(\mathbf{f} \given \mathbf{e})$ that captures
word and phrase correspondence (trained on parallel corpora), and a
\emph{language model} $P(\mathbf{e})$ that ensures fluent output (trained on
large monolingual corpora).  Decoding uses beam search over the combined
objective.

SMT was replaced by neural MT because: (1) NMT is end-to-end---a single
encoder-decoder model replaces the complex SMT pipeline of alignment models,
phrase tables, reordering models, and feature-weight tuning.  (2) NMT produces
more fluent output by modeling the full target-side context, not just n-gram
language model scores.  (3) NMT handles long-range reordering naturally via
attention, while SMT reordering models are local.  (4) Subword tokenization
in NMT handles morphologically rich languages without exploding phrase tables.

\redflags
\begin{itemize}
    \item Not knowing the Bayesian decomposition or why $P(\mathbf{f})$ drops
          out of the argmax.
    \item Saying ``SMT was bad'' without articulating specific limitations.
    \item Not knowing what a phrase table is or what alignment means.
\end{itemize}

\interviewtest{Historical knowledge of NLP, understanding of Bayesian
reasoning, and ability to articulate why paradigm shifts happen.}

\goodgreat
\begin{itemize}
    \item \badgeLfive{L5}: Explains the decomposition clearly, gives 2--3
          reasons for the shift to NMT.
    \item \badgeLsix{L6}: Discusses IBM alignment models, phrase-based SMT,
          and the log-linear model that combined features.  Explains
          BLEU score as the evaluation metric developed for SMT.
    \item \badgeLseven{L7}: Connects the noisy channel model to other
          applications (speech recognition, spelling correction), discusses
          how modern LLMs can be viewed as implicit language models replacing
          both components, and notes that the noisy channel idea has been
          revived in some neural approaches (channel models for domain
          adaptation).
\end{itemize}

\followups{What are IBM alignment models?  What is phrase-based SMT?  How
does the noisy channel model appear in speech recognition?  Has the noisy
channel idea been used with neural models?}
\end{interviewq}

\vspace{0.5em}

\begin{interviewq}
\textbf{Q: You are building a named entity recognizer for a new domain
(e.g., legal documents) with only 500 labeled sentences.  Walk me through
your approach.}

\badgeLfive{L5} \quad \badgetype{Applied} \quad \badgefreq{\freqmed}

\stronganswer
With only 500 labeled sentences, data efficiency is the primary concern.
My approach:

\textbf{Step 1: Start with a strong pretrained model.}  Use a domain-relevant
pretrained language model (e.g., Legal-BERT or a general BERT/RoBERTa model)
as the feature extractor.  Pretraining provides rich representations that
compensate for limited task-specific data.

\textbf{Step 2: Add a CRF layer.}  With limited data, the CRF transition
matrix is especially valuable because it enforces structural constraints on
label sequences that the model cannot learn from sparse examples.  A
BERT-CRF architecture should outperform BERT-softmax in this low-data regime.

\textbf{Step 3: Data augmentation.}  Apply entity-preserving augmentation:
synonym replacement for non-entity tokens, entity swapping (replace entity
mentions with other entities of the same type), and sentence-level paraphrasing
using an LLM.

\textbf{Step 4: Active learning.}  Use the initial model to identify
unlabeled sentences where it is most uncertain, and prioritize those for
annotation.  This maximizes the value of each labeled example.

\textbf{Step 5: Few-shot with LLMs as augmentation.}  Use GPT-4 or a similar
model with few-shot prompting to generate additional labeled examples.  These
silver-label examples can supplement the gold data, especially for common
entity types.

\textbf{Step 6: Evaluation.}  Use entity-level F1 (not token-level) as the
primary metric.  Split the 500 sentences into train/dev/test (e.g.,
350/75/75), and use cross-validation to get reliable estimates given the
small size.

\redflags
\begin{itemize}
    \item Proposing to train a model from scratch---with 500 examples, you
          must use transfer learning.
    \item Not mentioning the CRF layer, which is most valuable in low-data
          settings.
    \item Ignoring data augmentation or active learning strategies.
    \item Using accuracy instead of entity-level F1 for evaluation.
\end{itemize}

\interviewtest{Practical NLP engineering skills, ability to reason about
data efficiency, and knowledge of the full toolkit from pretrained models
to data augmentation to active learning.}

\goodgreat
\begin{itemize}
    \item \badgeLfive{L5}: Proposes fine-tuning a pretrained model with CRF,
          mentions data augmentation.
    \item \badgeLsix{L6}: Discusses specific augmentation strategies, active
          learning, few-shot LLM labeling, and evaluation methodology for
          small datasets (cross-validation, confidence intervals).
    \item \badgeLseven{L7}: Discusses continued pretraining on domain text
          before fine-tuning, curriculum learning, the tradeoff between model
          complexity and data availability, and how to set up a human-in-the-loop
          annotation pipeline for iterative improvement.
\end{itemize}

\followups{How do you handle class imbalance in NER?  What entity-level F1
metric do you use (strict vs.\ partial matching)?  How would you decide
between a CRF-based approach and an LLM prompting approach?}
\end{interviewq}
